\section{Several units of descent and color coding}\label{sec:multiunit}

The connected witness theorem answers one-unit descent. Requiring a decrease
by at least $k$ tetrahedra may require several independent regions. Their
locations need not be enumerated with a $t^k$ factor: their total original
footprint remains bounded, which permits standard deterministic color coding.
This section gives a further algorithmic theorem. The delivered code includes
the exact packing primitive for a supplied coloring and finite audits; it does
not implement the full deterministic perfect-hash-family construction.

\begin{theorem}[Multiunit localization]\label{thm:k-local}
Suppose a trace with at most $U$ upward events reduces $T$ by at least
$k\ge1$ tetrahedra. Then there is such a trace with $u\le U$, $d=u+k$,
length $2u+k$, and at most $k$ birth components. Each component has positive
loss and connected actual original footprint. Their disjoint union has size
at most $3U+3k$.
\end{theorem}
\begin{proof}
Choose a shortest trace meeting the threshold. It first reaches loss $k$
at its final event, so $d=u+k$. If a component had nonpositive loss, deleting
it using \cref{thm:projection} would preserve loss at least $k$ and shorten
the trace. All component losses $\ell_j$ are therefore positive integers,
with $\sum_j\ell_j=k$, so the number $c$ of components is at most $k$.
Their footprints are disjoint and connected by \cref{lem:connected}. Summing
\cref{lem:incidence} yields
\[
 |F|\le u+2d+c=3u+2k+c\le3U+3k.
\]
\end{proof}

\begin{lemma}[Composing disjoint original footprints]\label{lem:compose}
Let $\sigma_1,\ldots,\sigma_s$ be legal traces from the same marked source
whose actual original footprints are pairwise disjoint. They can be replayed
in any order, transporting the exterior attachments at each step. The
combined upward cost and tetrahedron loss are the sums of the individual
costs and losses.
\end{lemma}
\begin{proof}
Give births in different traces disjoint names. Every incarnation consumed
by one trace is either one of its original footprint cells or a birth
descended from those cells. Distinct traces therefore have disjoint consumed
incarnation sets after transport. At the common source their first events
are simultaneously legal and commute. Inductively construct the grid of
interleavings: at every square, \cref{lem:forward} transports the next event
of either trace across the next event of the other, retaining the same
relative-boundary attachments. This proves that both complete orders are
legal. Apply the two-trace argument repeatedly. Marking agreement follows
from the same local transport lemma, and the count assertions follow by
adding event counts.
\end{proof}

This lemma does not require an actual footprint to be a ball or an induced
subcomplex with embedded closure. Individual formal replacements, with their
boundary ports, are the units of the proof. It also does not claim that
arbitrary traces with overlapping original footprints compose.

\subsection{Enumerating candidate items}

Set $p=\min(t,3U+3k)$. Over every connected original cover of size at most
$p$, enumerate the full bounded-upward region family and retain every
positive-loss endpoint. Each item $j$ records
\[
 (A_j,a_j,b_j,\sigma_j),
\]
Here $A_j$ is its \emph{actual} consumed original footprint, $a_j\le U$
is its upward cost, and $\sigma_j$ is a replayable source-bound trace.
The positive profit is capped: $b_j=\min(k,\loss(\sigma_j))$. The total item
count and generation cost are bounded by
\begin{equation}\label{eq:items}
 M\le t\,2^{O(p+U)}\poly(t+U+B).
\end{equation}
This is the generic connected-region commitment bound. A particular actual
footprint can appear under several covers; retaining duplicate items does
not affect the upper bound or correctness.

The item generator must continue through intermediate descents. Calling a
first-descent routine and collecting only its first answer would not enumerate
the needed component alternatives. Birth-connectedness need not be tested
for every retained item: extra legal items remain sound, while every component
of the minimum witness in \cref{thm:k-local} is necessarily represented.

\subsection{Colorful packing with original witnesses}

A family $\mathcal H$ of functions $[t]\to[p]$ is $p$-perfect if every subset
of at most $p$ source indices is mapped injectively by at least one member.
The standard deterministic construction of Naor, Schulman, and Srinivasan
has size $e^p p^{O(\log p)}\log(t+2)$ and total construction time bounded by
$2^{O(p)}\poly(t)$ \cite{nss}. We need this total-time statement; we do not
assume a constant-delay hash enumerator. Color coding itself is classical
\cite{color-coding}; an explicit modern treatment of the stronger delay
requirements is available in \cite{gutin-color}.

For a fixed coloring $h\in\mathcal H$, discard an item if $h$ is not
injective on $A_j$. Otherwise store its nonempty color mask
$C_j=h(A_j)\subseteq[p]$. Use the following ordinary $0/1$ dynamic program:
\begin{align*}
 D_0(\varnothing,0)&=0,\qquad D_0(S,a)=-\infty\text{ otherwise},\\
 D_j(S,a)&\gets D_{j-1}(S,a),\\
 D_j(S\cup C_j,a+a_j)&\gets
 \max\{D_j(S\cup C_j,a+a_j),\min(k,D_{j-1}(S,a)+b_j)\}
\end{align*}
for every finite state with $S\cap C_j=\varnothing$ and $a+a_j\le U$.
The state value is the best capped total loss from the first $j$ items,
using exactly mask $S$ and upward cost $a$. Take/skip predecessor data
retains actual item identities for reconstruction; a linear combination of
traces is never treated as a geometric witness.

The number of states is $2^p(U+1)$, and a direct item scan costs
$O(M2^p(U+1))$. Positive items have nonempty footprints, so the empty-mask
degeneracy does not arise. Mask disjointness implies actual source-cell
disjointness, even though distinct source cells may share colors. By
\cref{lem:compose}, a successful collection is geometrically replayable and
has total upward cost at most $U$ and loss at least $k$.

\begin{theorem}[Fixed-parameter multiunit descent]\label{thm:k-fpt}
For generalized formal-bipyramid moves, existence of a trace with at most
$U$ total upward moves reducing a valid $t$-tetrahedron source by at least
$k$ tetrahedra is decidable deterministically in
\begin{equation}\label{eq:k-fpt}
 2^{O(U+k)}\poly(t+B)
\end{equation}
bit time. A positive result admits a source-bound geometric trace certificate.
\end{theorem}
\begin{proof}
Trivial impossible thresholds, including $k\ge t$ for a nonempty connected
source, may be answered directly. In the remaining cases construct the
items above and run the packing dynamic program for every coloring in
$\mathcal H$. Soundness follows from \cref{lem:compose}.

For completeness, take the minimum-length witness of \cref{thm:k-local}.
Its component footprints are individually connected, pairwise disjoint,
and their union has at most $p$ vertices. Each component trace occurs among
the items. Some $h\in\mathcal H$ colors the entire union injectively, so
these items have pairwise disjoint color masks. Their upward costs sum to
at most $U$ and their profits sum to $k$; the dynamic program finds a
successful collection. Enumerating every coloring therefore decides the
question completely.

Multiply \eqref{eq:items}, the $2^p(U+1)$ dynamic-program factor, and the
$2^{O(p)}\log(t+2)$ hash-family factor. Since $p\le3(U+k)$, the result is
\eqref{eq:k-fpt}, with polynomial source and replay costs included.
\end{proof}

The weaker direct alternative is to enumerate covers with up to $k$
components, which gives $t^k2^{O(U+k)}\poly(t+U+B)$. The color-coding
argument removes the $t^k$ factor. This is a consequence of the new bounded
total footprint and independent replay, composed with an established
derandomization theorem. It is not a claim that color coding is new.

With $U+k=O(\log t)$ the decision problem is polynomial; with
$U+k=O(\log^2 t)$ it is quasi-polynomial. The supplied-coloring prototype
tests the witness-preserving packing step, while \cref{thm:k-fpt} relies
on the cited complete perfect-family construction for deterministic
exhaustion over all possible locations.
